% Options for packages loaded elsewhere
\PassOptionsToPackage{unicode}{hyperref}
\PassOptionsToPackage{hyphens}{url}
\documentclass[
]{article}
\usepackage{xcolor}
\usepackage[margin=1in]{geometry}
\usepackage{amsmath,amssymb}
\setcounter{secnumdepth}{-\maxdimen} % remove section numbering
\usepackage{iftex}
\ifPDFTeX
  \usepackage[T1]{fontenc}
  \usepackage[utf8]{inputenc}
  \usepackage{textcomp} % provide euro and other symbols
\else % if luatex or xetex
  \usepackage{unicode-math} % this also loads fontspec
  \defaultfontfeatures{Scale=MatchLowercase}
  \defaultfontfeatures[\rmfamily]{Ligatures=TeX,Scale=1}
\fi
\usepackage{lmodern}
\ifPDFTeX\else
  % xetex/luatex font selection
\fi
% Use upquote if available, for straight quotes in verbatim environments
\IfFileExists{upquote.sty}{\usepackage{upquote}}{}
\IfFileExists{microtype.sty}{% use microtype if available
  \usepackage[]{microtype}
  \UseMicrotypeSet[protrusion]{basicmath} % disable protrusion for tt fonts
}{}
\makeatletter
\@ifundefined{KOMAClassName}{% if non-KOMA class
  \IfFileExists{parskip.sty}{%
    \usepackage{parskip}
  }{% else
    \setlength{\parindent}{0pt}
    \setlength{\parskip}{6pt plus 2pt minus 1pt}}
}{% if KOMA class
  \KOMAoptions{parskip=half}}
\makeatother
\usepackage{color}
\usepackage{fancyvrb}
\newcommand{\VerbBar}{|}
\newcommand{\VERB}{\Verb[commandchars=\\\{\}]}
\DefineVerbatimEnvironment{Highlighting}{Verbatim}{commandchars=\\\{\}}
% Add ',fontsize=\small' for more characters per line
\usepackage{framed}
\definecolor{shadecolor}{RGB}{248,248,248}
\newenvironment{Shaded}{\begin{snugshade}}{\end{snugshade}}
\newcommand{\AlertTok}[1]{\textcolor[rgb]{0.94,0.16,0.16}{#1}}
\newcommand{\AnnotationTok}[1]{\textcolor[rgb]{0.56,0.35,0.01}{\textbf{\textit{#1}}}}
\newcommand{\AttributeTok}[1]{\textcolor[rgb]{0.13,0.29,0.53}{#1}}
\newcommand{\BaseNTok}[1]{\textcolor[rgb]{0.00,0.00,0.81}{#1}}
\newcommand{\BuiltInTok}[1]{#1}
\newcommand{\CharTok}[1]{\textcolor[rgb]{0.31,0.60,0.02}{#1}}
\newcommand{\CommentTok}[1]{\textcolor[rgb]{0.56,0.35,0.01}{\textit{#1}}}
\newcommand{\CommentVarTok}[1]{\textcolor[rgb]{0.56,0.35,0.01}{\textbf{\textit{#1}}}}
\newcommand{\ConstantTok}[1]{\textcolor[rgb]{0.56,0.35,0.01}{#1}}
\newcommand{\ControlFlowTok}[1]{\textcolor[rgb]{0.13,0.29,0.53}{\textbf{#1}}}
\newcommand{\DataTypeTok}[1]{\textcolor[rgb]{0.13,0.29,0.53}{#1}}
\newcommand{\DecValTok}[1]{\textcolor[rgb]{0.00,0.00,0.81}{#1}}
\newcommand{\DocumentationTok}[1]{\textcolor[rgb]{0.56,0.35,0.01}{\textbf{\textit{#1}}}}
\newcommand{\ErrorTok}[1]{\textcolor[rgb]{0.64,0.00,0.00}{\textbf{#1}}}
\newcommand{\ExtensionTok}[1]{#1}
\newcommand{\FloatTok}[1]{\textcolor[rgb]{0.00,0.00,0.81}{#1}}
\newcommand{\FunctionTok}[1]{\textcolor[rgb]{0.13,0.29,0.53}{\textbf{#1}}}
\newcommand{\ImportTok}[1]{#1}
\newcommand{\InformationTok}[1]{\textcolor[rgb]{0.56,0.35,0.01}{\textbf{\textit{#1}}}}
\newcommand{\KeywordTok}[1]{\textcolor[rgb]{0.13,0.29,0.53}{\textbf{#1}}}
\newcommand{\NormalTok}[1]{#1}
\newcommand{\OperatorTok}[1]{\textcolor[rgb]{0.81,0.36,0.00}{\textbf{#1}}}
\newcommand{\OtherTok}[1]{\textcolor[rgb]{0.56,0.35,0.01}{#1}}
\newcommand{\PreprocessorTok}[1]{\textcolor[rgb]{0.56,0.35,0.01}{\textit{#1}}}
\newcommand{\RegionMarkerTok}[1]{#1}
\newcommand{\SpecialCharTok}[1]{\textcolor[rgb]{0.81,0.36,0.00}{\textbf{#1}}}
\newcommand{\SpecialStringTok}[1]{\textcolor[rgb]{0.31,0.60,0.02}{#1}}
\newcommand{\StringTok}[1]{\textcolor[rgb]{0.31,0.60,0.02}{#1}}
\newcommand{\VariableTok}[1]{\textcolor[rgb]{0.00,0.00,0.00}{#1}}
\newcommand{\VerbatimStringTok}[1]{\textcolor[rgb]{0.31,0.60,0.02}{#1}}
\newcommand{\WarningTok}[1]{\textcolor[rgb]{0.56,0.35,0.01}{\textbf{\textit{#1}}}}
\usepackage{graphicx}
\makeatletter
\newsavebox\pandoc@box
\newcommand*\pandocbounded[1]{% scales image to fit in text height/width
  \sbox\pandoc@box{#1}%
  \Gscale@div\@tempa{\textheight}{\dimexpr\ht\pandoc@box+\dp\pandoc@box\relax}%
  \Gscale@div\@tempb{\linewidth}{\wd\pandoc@box}%
  \ifdim\@tempb\p@<\@tempa\p@\let\@tempa\@tempb\fi% select the smaller of both
  \ifdim\@tempa\p@<\p@\scalebox{\@tempa}{\usebox\pandoc@box}%
  \else\usebox{\pandoc@box}%
  \fi%
}
% Set default figure placement to htbp
\def\fps@figure{htbp}
\makeatother
\setlength{\emergencystretch}{3em} % prevent overfull lines
\providecommand{\tightlist}{%
  \setlength{\itemsep}{0pt}\setlength{\parskip}{0pt}}
\usepackage{bookmark}
\IfFileExists{xurl.sty}{\usepackage{xurl}}{} % add URL line breaks if available
\urlstyle{same}
\hypersetup{
  pdftitle={Bayesian Logistic Regression},
  hidelinks,
  pdfcreator={LaTeX via pandoc}}

\title{Bayesian Logistic Regression}
\author{}
\date{\vspace{-2.5em}2026-04-17}

\begin{document}
\maketitle

\subsection{Introduction}\label{introduction}

Bayesian logistic regression keeps the same Bernoulli likelihood and
logit link as the maximum-likelihood version but replaces the point
estimate of each coefficient with a full posterior distribution. The
biggest practical pay-off is robustness in small samples or in the
presence of complete or quasi-complete separation, where MLE drifts to
\(\pm\infty\) but a weakly informative prior pulls the estimate to a
finite, defensible value. The posterior also feeds directly into
propagated uncertainty for predicted probabilities, an awkward
computation in the frequentist world.

\subsection{Prerequisites}\label{prerequisites}

Familiarity with classical logistic regression and odds ratios, an idea
of what a prior and a posterior are, and access to a Stan-capable
backend (\texttt{rstan} or \texttt{cmdstanr}).

\subsection{Theory}\label{theory}

The model is

\[y_i \sim \mathrm{Bernoulli}\bigl(\mathrm{logit}^{-1}(X_i \beta)\bigr),\]

with priors \(\beta \sim \mathcal N(0, \sigma_\beta^2)\) on the
regression coefficients and a separate, often wider prior on the
intercept because the average log-odds of success can lie outside the
typical predictor scale. A \(\mathcal N(0, 2.5)\) prior on standardised
slopes is the default in \texttt{rstanarm} and a sensible starting point
in \texttt{brms}. Under separation the likelihood becomes monotone in
\(\beta\), and any proper prior --- even a very wide one --- yields a
well-defined posterior with finite credible intervals; this is the
Bayesian analogue of Firth's penalised likelihood.

\subsection{Assumptions}\label{assumptions}

Independent Bernoulli observations conditional on covariates, a logit
link that is a reasonable approximation to the true response surface,
and valid (proper) priors. As in MLE logistic regression, the linear
predictor should not contain redundant columns.

\subsection{R Implementation}\label{r-implementation}

\begin{Shaded}
\begin{Highlighting}[]
\FunctionTok{library}\NormalTok{(brms)}

\FunctionTok{set.seed}\NormalTok{(}\DecValTok{2026}\NormalTok{)}
\NormalTok{d }\OtherTok{\textless{}{-}} \FunctionTok{data.frame}\NormalTok{(}\AttributeTok{x =} \FunctionTok{rnorm}\NormalTok{(}\DecValTok{200}\NormalTok{))}
\NormalTok{d}\SpecialCharTok{$}\NormalTok{y }\OtherTok{\textless{}{-}} \FunctionTok{rbinom}\NormalTok{(}\DecValTok{200}\NormalTok{, }\DecValTok{1}\NormalTok{, }\FunctionTok{plogis}\NormalTok{(}\SpecialCharTok{{-}}\DecValTok{1} \SpecialCharTok{+} \FloatTok{1.2} \SpecialCharTok{*}\NormalTok{ d}\SpecialCharTok{$}\NormalTok{x))}

\NormalTok{fit }\OtherTok{\textless{}{-}} \FunctionTok{brm}\NormalTok{(y }\SpecialCharTok{\textasciitilde{}}\NormalTok{ x, }\AttributeTok{data =}\NormalTok{ d, }\AttributeTok{family =}\NormalTok{ bernoulli,}
           \AttributeTok{prior =} \FunctionTok{prior}\NormalTok{(}\FunctionTok{normal}\NormalTok{(}\DecValTok{0}\NormalTok{, }\FloatTok{2.5}\NormalTok{), }\AttributeTok{class =} \StringTok{"b"}\NormalTok{) }\SpecialCharTok{+}
                   \FunctionTok{prior}\NormalTok{(}\FunctionTok{normal}\NormalTok{(}\DecValTok{0}\NormalTok{, }\DecValTok{5}\NormalTok{), }\AttributeTok{class =} \StringTok{"Intercept"}\NormalTok{),}
           \AttributeTok{chains =} \DecValTok{4}\NormalTok{, }\AttributeTok{iter =} \DecValTok{2000}\NormalTok{, }\AttributeTok{refresh =} \DecValTok{0}\NormalTok{)}
\FunctionTok{summary}\NormalTok{(fit)}

\CommentTok{\# Odds ratios}
\FunctionTok{exp}\NormalTok{(}\FunctionTok{fixef}\NormalTok{(fit))}
\end{Highlighting}
\end{Shaded}

\subsection{Output \& Results}\label{output-results}

\texttt{summary(fit)} reports posterior mean, standard error, and 95 \%
CrI on the log-odds scale together with \(\hat R\) and effective sample
size; \texttt{exp(fixef(fit))} returns the same summaries on the
odds-ratio scale. \texttt{posterior\_predict()} delivers Bernoulli
replicates that feed \texttt{pp\_check()} to test calibration.

\subsection{Interpretation}\label{interpretation}

A reporting sentence: ``A Bayesian logistic regression with weakly
informative \(\mathcal N(0, 2.5)\) priors estimated the odds ratio for
\(x\) at 3.4 (95 \% CrI 2.5 to 4.5); under the model, a
one-standard-deviation increase in \(x\) triples the odds of success.''
When predicted probabilities are required, summarise
\texttt{posterior\_epred()} rather than transforming \texttt{fixef()} by
hand --- only the former propagates the joint uncertainty in intercept
and slope.

\subsection{Practical Tips}\label{practical-tips}

\begin{itemize}
\tightlist
\item
  For binary regression with rare outcomes, prefer a
  \(\mathcal N(0, 1)\) prior on standardised slopes; the default
  \(\mathcal N(0, 2.5)\) allows odds ratios up to about
  \(\exp(5) \approx 150\), which can be implausibly extreme.
\item
  Compare predictive performance with \texttt{loo()} rather than
  likelihood-ratio tests; LOO is well-defined under proper priors and
  gives Pareto-\(k\) diagnostics for influential points.
\item
  Hierarchical extensions are a one-line addition:
  \texttt{y\ \textasciitilde{}\ x\ +\ (1\ \textbar{}\ site)} lets
  baseline log-odds vary across sites and shrinks small-cluster
  estimates toward the population mean.
\item
  For severe separation, fit with
  \texttt{prior(normal(0,\ 1.5),\ class\ =\ "b")} rather than ML's Firth
  correction; the Bayesian solution is a single line of code and
  propagates uncertainty consistently.
\item
  Always inspect \texttt{pp\_check(fit,\ type\ =\ "bars")} for binary
  outcomes; deviations between observed and replicated proportions
  reveal calibration problems that summary tables hide.
\end{itemize}

\subsection{R Packages Used}\label{r-packages-used}

\texttt{brms} for model specification, \texttt{bayesplot} for posterior
predictive plots, \texttt{loo} for cross-validation and model
comparison, and \texttt{tidybayes} if you want tidy tibbles of posterior
draws for downstream plotting.

\subsection{For Reviewers}\label{for-reviewers}

\textbf{What to look for in a paper using this method.}

\begin{itemize}
\tightlist
\item
  \textbf{Common misapplications.}
\item
  \textbf{Diagnostics that should be reported but often aren't.}
\item
  \textbf{Red flags in tables and figures.}
\item
  \textbf{What to verify.}
\item
  \textbf{What an adequate Methods paragraph must contain.}
\end{itemize}

\subsection{See also --- labs in this
chapter}\label{see-also-labs-in-this-chapter}

\begin{itemize}
\tightlist
\item
  \href{labs/lab-bayesian-thinking-control-vs-decision-error.Rmd}{Bayesian
  thinking; α control vs decision error}
\item
  \href{labs/lab-brms-stan-loo-hierarchical-models.Rmd}{brms / Stan,
  LOO, hierarchical models}
\end{itemize}

\end{document}
